\documentclass[11pt]{article}
\usepackage[T1]{fontenc}
\usepackage{textcomp}
\usepackage[margin=0.75in]{geometry}
\usepackage{amsmath,amssymb,amsthm}
\usepackage{array,booktabs}
\usepackage{hyperref}
\setlength{\parindent}{0pt}
\newtheorem{theorem}{Theorem}
\theoremstyle{definition}
\newtheorem{definition}{Definition}
\title{Flow Proof Artifact: Ratio}
\author{Generated by \texttt{flow doc proof}}
\date{Flow Proof Book}
\begin{document}
\maketitle
Eudoxus proportion definitions on comparable magnitudes.\\[0.5em]
\textbf{Source.} euclid — Elements, Book V, Definition 5\\[0.5em]
\bigskip
\noindent{\large \textbf{Definition 1.} \textit{A magnitude is in the same ratio to a second as a third to a fourth when equimultiples match} \hfill Ratio \cdot proportion \cdot same ratio means matching equimultiples}\par
\smallskip\noindent\textit{Coordinate: Ratio \cdot proportion \cdot same ratio means matching equimultiples}\par
\smallskip\noindent\textit{Source: euclid}\par
\medskip
\noindent\textbf{Goal.} A magnitude is in the same ratio to a second as a third to a fourth when equimultiples match\par
\noindent$\displaystyle \forall a \in Magnitude \forall b \in Magnitude \forall c \in Magnitude \forall d \in Magnitude\quad ratio(a, b) = ratio(c, d)$\par
\medskip
\noindent
\begin{tabular}{@{} >{\raggedright\arraybackslash}p{0.06\textwidth} >{\raggedright\arraybackslash}p{0.47\textwidth} >{\raggedleft\arraybackslash}p{0.41\textwidth} @{}}
\textbf{\#} & \textbf{Proof} & \textbf{Mathematics} \\
\midrule
\textbf{1.} & We stipulate same ratio means matching equimultiples for proportion on Ratio — this is a definition, not a derived fact. & \\[0.45em]
\textbf{2.} & This follows directly from the definition: ratio(a, b) equals ratio(c, d). Hence proven. & $\displaystyle ratio(a, b) = ratio(c, d)$ \\[0.45em]
\bottomrule
\end{tabular}
\label{thm:Ratio-proportion-same_ratio_means_matching_equimultiples}
\par\medskip\hrule
\bigskip
\noindent{\large \textbf{Derived fact 2.} \textit{Every magnitude is in the same ratio to itself} \hfill Ratio \cdot proportion \cdot every magnitude is proportional to itself}\par
\smallskip\noindent\textit{Coordinate: Ratio \cdot proportion \cdot every magnitude is proportional to itself}\par
\smallskip\noindent\textit{Source: euclid}\par
\medskip
\noindent\textbf{Goal.} Every magnitude is in the same ratio to itself\par
\noindent$\displaystyle \forall a \in Magnitude\quad ratio(a, a) = ratio(a, a)$\par
\medskip
\noindent
\begin{tabular}{@{} >{\raggedright\arraybackslash}p{0.06\textwidth} >{\raggedright\arraybackslash}p{0.47\textwidth} >{\raggedleft\arraybackslash}p{0.41\textwidth} @{}}
\textbf{\#} & \textbf{Proof} & \textbf{Mathematics} \\
\midrule
\textbf{1.} & We prove that every magnitude is proportional to itself for proportion on Ratio. & \\[0.45em]
\textbf{2.} & We can deduce that ratio(a, a) equals ratio(a, a). Hence proven. & $\displaystyle ratio(a, a) = ratio(a, a)$ \\[0.45em]
\bottomrule
\end{tabular}
\label{thm:Ratio-proportion-every_magnitude_is_proportional_to_itself}
\par\medskip\hrule
\end{document}
